%!TEX program = uplatex
\RequirePackage{plautopatch}

\documentclass[uplatex,dvipdfmx]{jlreq}

% 余白調整（1ページ収め）
\usepackage[top=18mm,bottom=20mm,left=20mm,right=20mm]{geometry}

\usepackage{amsmath,amssymb}
\usepackage{url}

\pagestyle{empty}

% 行間少し詰める
\renewcommand{\baselinestretch}{0.96}

%--------------------------------
% 講演マクロ
%--------------------------------

\newcommand{\talk}[3]{%
\noindent
#1\quad #2\par
\hspace*{2em}#3\par\smallskip
}

%--------------------------------
% タイトル情報
%--------------------------------

\newcommand{\EwMnumber}{第10回}
\newcommand{\EwMtitle}{応用特異点論\\[2mm]\large 数学（厳密科学）としてのカタストロフ理論をめざして}
\newcommand{\EwMdate}{1999年2月5日（金）14:30 ～ 2月6日（土）17:00}
\newcommand{\EwMplace}{東京都文京区春日1-13-27 中央大学理工学部5号館}

%--------------------------------
% タイトル表示
%--------------------------------

\newcommand{\EwMtitleblock}{

\vspace*{-5mm}

\begin{center}

{\Large ENCOUNTER with MATHEMATICS}

\vspace{2mm}

{\large 数学との遭遇}

\vspace{4mm}

{\Large \bfseries \EwMnumber\ \EwMtitle}

\vspace{4mm}

\EwMdate

\vspace{2mm}

於：\EwMplace

\end{center}

\vspace{2mm}

}

\begin{document}

\EwMtitleblock

%--------------------------------
% プログラム
%--------------------------------

\section*{プログラム}

\subsection*{1日目（2月5日（金））}

\talk
{14:30～15:50}
{応用特異点論概説}
{泉屋 周一 氏（北大・理）}

\talk
{16:30～17:30}
{応用特異点論の基礎 I}
{石川 剛郎 氏（北大・理）}

\subsection*{2日目（2月6日（土））}

\talk
{10:30～11:30}
{応用特異点論の基礎 II}
{石川 剛郎 氏（北大・理）}

\talk
{11:50～12:50}
{一階偏微分方程式への応用}
{泉屋 周一 氏（北大・理）}

\talk
{14:40～15:40}
{微分幾何学への応用}
{泉屋 周一 氏（北大・理）}

\talk
{16:00～17:00}
{微分位相幾何学への応用}
{佐伯 修 氏（広島大・理）}

%--------------------------------
% 案内文
%--------------------------------

\bigskip

別紙の趣旨に沿った集会の第10回を以上のような予定で開催いたします。  
非専門家向けに入門的な講演をお願い致しました。  
多くの方々のご参加をお待ちしております。  
講演者による講演内容へのご案内を別紙に添付いたしますのでご覧下さい。

%--------------------------------
% 連絡先
%--------------------------------

\section*{連絡先}

112 東京都文京区春日1-13-27 中央大学理工学部数学教室  

tel : 03-3817-1745  

ENCOUNTER with MATHEMATICS : e-mail : encmath@math.chuo-u.ac.jp  
homepage : \url{http://www.math.chuo-u.ac.jp/ENCwMATH}  

三松 佳彦 : yoshi@math.chuo-u.ac.jp / 高倉 樹 : takakura@math.chuo-u.ac.jp

\end{document}
